% !TeX program = pdflatex
\ifdefined\pdfoutput\pdfoutput=1\fi % arXiv builds with pdfLaTeX when it sees this line
\documentclass[titlestyle=card]{udemmila}
% titlestyle = card | band | rule          (UdeM and Mila logos)
%            | labcard | labband | masthead (plus the lab wordmark, no edge bars)
%            | cover | hero | teaser | split | swiss | highlights
%            | letterhead | bilingual | motif | minimal   (see Appendix B)

% The class already loads graphicx, booktabs, amsmath, amssymb, xcolor
% (with the table option), tcolorbox, hyperref, cleveref and natbib.
\usepackage{multirow}
\usepackage{pgfplots}
\pgfplotsset{compat=1.18}

% Front matter %%%%%%%%%%%%%%%%%%%%%%%%%%%%%%%%%%%%%%%%%%%%%%%%%%%%%%%%%%%%%%%%%%

\title{Title of Your Paper in Two Lines at Most:\\ A Subtitle That Names the Method}
\runningtitle{Short Title for the Running Head}
\reportlabel{Technical Report}
\lab{Bang}{Lab}                                        % wordmark, used by the lab styles
\labdescriptor{Universit\'e de Montr\'eal \textperiodcentered{} Mila}

\author[1,2,*]{First Author}
\author[1,2,*]{Second Author}
\author[1,3]{Third Author}
\author[1,2,3,\dagger]{Senior Author}

\affiliation[1]{Mila -- Quebec AI Institute}
\affiliation[2]{Universit\'e de Montr\'eal}
\affiliation[3]{Canada CIFAR AI Chair}
\contribution[*]{Equal contribution}
\contribution[\dagger]{Corresponding author}

\abstract{This template typesets arXiv preprints for authors at Universit\'e de
Montr\'eal and Mila. It keeps what works in the industry preprint styles of Meta
FAIR, ByteDance Seed, ServiceNow and LeapLab, a front matter that reads as one
block, sans-serif headings over a serif body, and metadata a reader can click,
and sets them in Mila purple and UdeM blue with the BangLab wordmark. A single
class option chooses the title page. The author, affiliation and metadata commands follow the Meta
FAIR convention, so a paper written for the Meta FAIR, ServiceNow or ByteDance
Seed class moves here by changing its \texttt{\textbackslash documentclass}
line. Replace this paragraph with your abstract and keep it to one paragraph of
about 150 to 250 words.}

\date{\today}
\correspondence{\email{first.author@mila.quebec}, \email{senior.author@umontreal.ca}}
\github{\href{https://github.com/your-org/your-repo}{github.com/your-org/your-repo}}
\huggingface{\href{https://huggingface.co/your-org}{huggingface.co/your-org}}

% Printed only by the styles built around them (Appendix B): the teaser figure
% by teaser, key numbers by highlights, the report number and copyright line
% by letterhead and minimal, the French abstract by bilingual.
\teaser{%
  \begin{tikzpicture}
    \begin{axis}[
        width=\linewidth, height=4.4cm, ybar=2pt, bar width=10pt,
        ymin=0, ymax=80, ytick={0,20,40,60,80}, enlarge x limits=0.12,
        /pgf/number format/assume math mode=true,
        symbolic x coords={Bench-A,Bench-B,Bench-C,Bench-D,Average}, xtick=data,
        axis lines*=left, axis line style={draw=umrule}, tick style={draw=none},
        ymajorgrids, grid style={draw=black!8},
        ticklabel style={font=\sffamily\scriptsize, text=umgrey},
        ylabel={Score}, ylabel style={font=\sffamily\scriptsize, text=umgrey},
        legend style={font=\sffamily\scriptsize, draw=none, fill=none,
          at={(0.5,1.02)}, anchor=south, legend columns=2,
          /tikz/every even column/.append style={column sep=12pt}},
        area legend,
        nodes near coords style={font=\sffamily\tiny, text=umink},
        every node near coord/.append style={
          /pgf/number format/.cd, fixed, fixed zerofill, precision=1, assume math mode=true}]
      \addplot[fill=chartblue, draw=none]
        coordinates {(Bench-A,49.5) (Bench-B,60.2) (Bench-C,37.1) (Bench-D,66.4) (Average,53.3)};
      \addplot[fill=chartpurple, draw=none, nodes near coords]
        coordinates {(Bench-A,53.1) (Bench-B,63.7) (Bench-C,41.8) (Bench-D,68.0) (Average,56.7)};
      \legend{Strongest published baseline, Ours}
    \end{axis}
  \end{tikzpicture}}{Teaser figure placeholder, drawn from the placeholder scores of
  \cref{tab:results}. Replace it with the headline result of your paper.}
\highlight{56.7}{average score over the four benchmarks}
\highlight{+3.4}{points over the strongest published baseline}
\highlight{41.8}{on Bench-C, the hardest of the four}
\reportnumber{BangLab-TR-2026-01}
\copyrightnote{\textcopyright\ 2026 BangLab, Universit\'e de Montr\'eal and Mila.}
\frenchabstract{Ce gabarit met en page les prépublications arXiv des auteurs de
l'Université de Montréal et de Mila. Il reprend ce qui fonctionne dans les styles
de prépublication industriels de Meta FAIR, ByteDance Seed, ServiceNow et LeapLab,
une page de titre qui se lit d'un seul bloc, des titres sans empattement sur un
texte à empattement et des métadonnées cliquables, et les compose en violet Mila
et en bleu UdeM avec le nom de BangLab. Une seule option de classe choisit la page
de titre. Les commandes d'auteur, d'affiliation et de métadonnées suivent la
convention de Meta FAIR, si bien qu'un article écrit pour la classe de Meta FAIR,
de ServiceNow ou de ByteDance Seed passe à celle-ci en changeant sa ligne
\texttt{\textbackslash documentclass}. Remplacez ce paragraphe par votre résumé,
en un seul paragraphe de 150 à 250 mots.}

\begin{document}

\maketitle

\section{Introduction}
\label{sec:intro}

Industry labs publish their arXiv preprints in house styles that make a paper
recognisable before its title is read. Meta FAIR sets the whole front matter in
a tinted card \citep{hao2024coconut}, ByteDance Seed rules off a centred title
under a logo header \citep{sheng2026rlcer}, ServiceNow frames the card in its
brand green and adds link lines for code and models \citep{wang2026mempi}, and
LeapLab puts the abstract on a grey panel below a logo and date
\citep{yue2025limit}. This class brings those choices to Universit\'e de
Montr\'eal and Mila, with the official UdeM and Mila marks and colours.

The body text is Latin Modern, so mathematics such as the attention operator of
\citet{vaswani2017attention},
\begin{equation}
  \operatorname{Attention}(Q, K, V) = \operatorname{softmax}\!\left(\frac{QK^\top}{\sqrt{d_k}}\right) V,
  \label{eq:attention}
\end{equation}
matches the surrounding text. Everything that carries the brand, the title
block, headings, caption labels and the running head, is set in Inter, the
closest free match to Mila's own typeface. \Cref{fig:overview} is where the
overview figure of a paper would go, and \cref{sec:usage} explains the class
options and the front matter commands.

\begin{figure*}[t]
  \centering
  \begin{tikzpicture}[
      font=\sffamily\small,
      stage/.style={rounded corners=5pt, minimum height=1.25cm, minimum width=3.3cm,
        align=center, draw=milapurple!55, line width=0.6pt, fill=umtint},
      core/.style={stage, fill=milapurple, draw=milapurple, text=white},
      flow/.style={-latex, line width=0.8pt, draw=milapurple!70}]
    \node[stage] (in)  at (0,0)   {\textbf{Input}\\[1pt]\footnotesize task and context};
    \node[core]  (m)   at (4.6,0) {\textbf{Your method}\\[1pt]\footnotesize the component that is new};
    \node[stage] (out) at (9.2,0) {\textbf{Output}\\[1pt]\footnotesize answer and trace};
    \node[stage, draw=udemblue!60, fill=udemblue!7, minimum width=4.2cm]
      (fb) at (4.6,-1.85) {\textbf{Feedback}\\[1pt]\footnotesize reward, critique or memory};
    \draw[flow] (in) -- (m);
    \draw[flow] (m) -- (out);
    \draw[flow, draw=udemblue!70] (out.south) |- (fb.east);
    \draw[flow, draw=udemblue!70] (fb.north) -- (m.south);
  \end{tikzpicture}
  \caption{Overview figure placeholder. The method node is filled in Mila purple
  and the feedback path is drawn in UdeM blue, the two brand colours this class
  defines as \texttt{milapurple} and \texttt{udemblue}.}
  \label{fig:overview}
\end{figure*}

\section{Using the class}
\label{sec:usage}

\subsection{Class options}

The class builds on \texttt{article}, so every \texttt{article} option such as
\texttt{11pt} or \texttt{twocolumn} passes through unchanged. \Cref{tab:options}
lists the two options the class adds. The title style is the only choice that
changes how the paper looks at a glance. The \texttt{card} style follows Meta
FAIR and ServiceNow, \texttt{band} puts a full-width purple band with white
logos across the top of the first page, and \texttt{rule} follows ByteDance Seed
and LeapLab with colour logos in the header and a ruled, centred title. The
three lab styles add the \labname{} wordmark and draw no coloured edge bars.
\texttt{labcard} is the card with the wordmark in place of its edge,
\texttt{labband} sets the wordmark in white on a solid purple band, and
\texttt{masthead} opens the page with the wordmark and logos above a double
rule and sets the abstract as plain text. Ten further styles, from a full
cover page to a two-language title page, are described in \cref{app:styles}.

\begin{table*}[t]
  \centering
  \caption{Options added by \texttt{udemmila.cls}. Any other option is passed to
  \texttt{article}.}
  \label{tab:options}
  \small
  \begin{tabular}{@{}lll@{}}
    \toprule
    \textbf{Option} & \textbf{Values} & \textbf{Effect} \\
    \midrule
    \texttt{titlestyle} & \texttt{card} (default), \texttt{band}, \texttt{rule} & Layout of the first page \\
                        & \texttt{labcard}, \texttt{labband}, \texttt{masthead} & The same with the lab wordmark \\
                        & ten more, listed in \cref{tab:styles} & See \cref{app:styles} \\
    \texttt{numbers}    & flag & Numeric citations [1] instead of (Author, Year) \\
    \bottomrule
  \end{tabular}
\end{table*}

\subsection{Front matter}

The front matter is declared in the preamble and printed by
\texttt{\textbackslash maketitle}. \Cref{tab:frontmatter} lists the commands.
Authors take their affiliation marks as an optional argument, and the abstract
is a command rather than an environment, both as in the Meta FAIR class. Every
metadata line is optional, and \texttt{\textbackslash metadata} adds a line of
any other kind with a Font Awesome icon.

\begin{table*}[t]
  \centering
  \caption{Front matter commands, all used in the preamble.}
  \label{tab:frontmatter}
  \small
  \begin{tabular}{@{}ll@{}}
    \toprule
    \textbf{Command} & \textbf{Prints} \\
    \midrule
    \verb|\title{...}|                      & Title, with \verb|\\| for a line break \\
    \verb|\author[1,2,*]{Name}|             & Author with affiliation and contribution marks \\
    \verb|\affiliation[1]{Mila}|            & Affiliation for mark 1 \\
    \verb|\contribution[*]{Equal contribution}| & Note for mark $*$ \\
    \verb|\abstract{...}|                   & Abstract \\
    \verb|\date|, \verb|\correspondence|, \verb|\github|, \verb|\huggingface|, \verb|\projectpage| & Metadata line with its icon \\
    \verb|\metadata[Key][\faBook]{Value}|   & Any other metadata line \\
    \verb|\runningtitle{...}|               & Header text on every page after the first \\
    \verb|\reportlabel{...}|                & Small label on the title page, Preprint by default \\
    \verb|\lab{Bang}{Lab}|, \verb|\labdescriptor{...}| & Lab wordmark and the line under it (lab styles) \\
    \verb|\teaser{figure}{caption}|, \verb|\highlight{56.7}{label}| & Page-one figure (teaser), key numbers (highlights) \\
    \verb|\reportnumber{...}|, \verb|\copyrightnote{...}|, \verb|\frenchabstract{...}| & Used by letterhead, minimal and bilingual \\
    \bottomrule
  \end{tabular}
\end{table*}

\section{Figures, tables and boxes}
\label{sec:floats}

Captions open with a purple label in Inter. A caption that fits on one line is
centred, and a longer one is justified. Table captions sit above the table.
\Cref{tab:results} shows the recommended results layout, with booktabs rules, the
best entry in bold, the runner-up underlined and the proposed method on a row
tinted with \texttt{umtint}.

\begin{table*}[t]
  \centering
  \caption{Example results table. All numbers are placeholders.}
  \label{tab:results}
  \small
  \begin{tabular}{@{}lccccc@{}}
    \toprule
    \textbf{Method} & \textbf{Bench-A} & \textbf{Bench-B} & \textbf{Bench-C} & \textbf{Bench-D} & \textbf{Avg.} \\
    \midrule
    Base model            & 41.2 & 55.0 & 30.4 & 62.1 & 47.2 \\
    Prompting baseline    & 44.8 & 57.3 & 33.9 & 63.0 & 49.8 \\
    Published method      & \underline{49.5} & \underline{60.2} & \underline{37.1} & \underline{66.4} & \underline{53.3} \\
    \rowcolor{umtint}
    Ours                  & \textbf{53.1} & \textbf{63.7} & \textbf{41.8} & \textbf{68.0} & \textbf{56.7} \\
    \bottomrule
  \end{tabular}
\end{table*}

Two boxes cover the content arXiv papers in this area most often set apart.
The \texttt{takeaway} box states a finding the reader should keep, and the
\texttt{promptbox} holds a prompt or a model output in Inconsolata, with
verbatim content allowed inside.

\begin{takeaway}[Takeaway]
State the finding in one or two sentences with the number that supports it, for
example that the method improves the average score by 3.4 points over the
strongest published baseline in \cref{tab:results}.
\end{takeaway}

\begin{promptbox}[Prompt template for the judge model]
You are grading a response to the task below.\\
Task: \{task\}\\
Response: \{response\}\\
Reply with a score from 1 to 10 and one sentence of justification.
\end{promptbox}

\section{Conclusion}

Replace this sample text with your paper. The appendix below starts with
\texttt{\textbackslash beginappendix}, which opens a new page, prints the
appendix heading and lists the appendix sections, and
\texttt{\textbackslash beginappendix*} leaves the list out.

\bibliographystyle{plainnat}
\bibliography{references}

\beginappendix

\section{Porting a paper from another template}
\label{app:porting}

A paper written for the Meta FAIR, ServiceNow or ByteDance Seed class compiles
here after its \texttt{\textbackslash documentclass} line changes to
\texttt{udemmila}. The commands those classes define, including
\texttt{\textbackslash checkdata}, \texttt{\textbackslash morelinks},
\texttt{\textbackslash beginappendix} and \texttt{\textbackslash nm}, keep their
meaning, and asset paths such as their logo files are no longer needed. Author
marks may be written as text or as maths symbols, since
\texttt{\textbackslash dagger} and \texttt{\textbackslash ddagger} are mapped to
their text forms inside the marks.

\section{Title styles}
\label{app:styles}

\Cref{tab:styles} lists every value of \texttt{titlestyle}. All but the first
three carry the lab wordmark, put it in the running head, and draw the takeaway
box without a coloured edge.

\begin{table*}[t]
  \centering
  \caption{The sixteen title styles. The last column names the preamble commands
  a style prints beyond the common front matter.}
  \label{tab:styles}
  \small
  \begin{tabular}{@{}lp{10.4cm}l@{}}
    \toprule
    \textbf{Style} & \textbf{First page} & \textbf{Extra} \\
    \midrule
    \texttt{card}       & Tinted card with a purple-to-blue edge, logos beside the metadata & \\
    \texttt{band}       & Gradient purple band with white logos, abstract with a left rule & \\
    \texttt{rule}       & Colour logos in the header, centred title between purple rules & \\
    \texttt{labcard}    & The card without its edge, opened by the wordmark & \\
    \texttt{labband}    & Solid purple band with the white wordmark, abstract panel & \\
    \texttt{masthead}   & Wordmark and logos above a double rule, no boxes & \\
    \texttt{cover}      & A full purple cover page, the abstract on page two & \\
    \texttt{hero}       & A deep aubergine block from the top edge holding the title & \\
    \texttt{teaser}     & Title between a heavy and a light rule, teaser figure below the abstract & \verb|\teaser| \\
    \texttt{split}      & Authors and links in a left column, title and abstract on the right & \\
    \texttt{swiss}      & Oversized title, three-column grid, the whole paper in Inter & \\
    \texttt{highlights} & The abstract beside a panel of key numbers & \verb|\highlight| \\
    \texttt{letterhead} & A report series line, small-capitals title, serif headings & \verb|\reportnumber| \\
    \texttt{bilingual}  & French and English label, the abstract beside its French résumé & \verb|\frenchabstract| \\
    \texttt{motif}      & A pale network pattern after the Mila glyph in the top corner & \\
    \texttt{minimal}    & Logos and date in the header, run-in abstract, copyright line & \verb|\copyrightnote| \\
    \bottomrule
  \end{tabular}
\end{table*}

\section{Brand assets}
\label{app:brand}

The logos in \texttt{assets/} are vector PDFs converted from the SVG files
served by \url{https://mila.quebec}. The Mila wordmark keeps its own purple,
\texttt{\#662E7D}, and the UdeM signature is set in UdeM blue,
\texttt{\#0057AC}, the colour \url{https://www.umontreal.ca} uses. White versions
of both are used on the purple band.

\end{document}
